\section{组合计数与递推关系}

\subsection{基本计数原理}

\subsubsection{加法原理与乘法原理}

\begin{enumerate}
	\item \textbf{加法原理}：若完成一件事有 \(k\) 类互不重叠的方法，第 \(i\) 类有 \(n_i\) 种，则共有
	\[
	n_1+n_2+\cdots+n_k
	\]
	种方法。
	\item \textbf{乘法原理}：若完成一件事分 \(k\) 步，第 \(i\) 步有 \(n_i\) 种选择，则共有
	\[
	n_1n_2\cdots n_k
	\]
	种方法。
\end{enumerate}

\subsubsection{排列与组合}

\begin{enumerate}
	\item 不重复排列：
	\[
	P(n,r)=A_n^r=\frac{n!}{(n-r)!}.
	\]
	\item 不重复组合：
	\[
	C(n,r)=\binom nr=\frac{n!}{r!(n-r)!}.
	\]
	\item 圆排列：
	\[
	(n-1)!.
	\]
	\item 含重复元素的全排列：若 \(n\) 个元素中有 \(n_1,n_2,\dots,n_k\) 个分别相同，且 \(\sum n_i=n\)，则不同排列数为
	\[
	\frac{n!}{n_1!n_2!\cdots n_k!}.
	\]
	\item 允许重复的 \(r\) 排列：
	\[
	n^r.
	\]
\end{enumerate}

\begin{question}{例题1}
	由数字 \(1,2,3,4,5\) 组成没有重复数字的三位数，共有多少个？
	
	\begin{solution}
	这是从 \(5\) 个数字中选 \(3\) 个并排列：
	\[
	P(5,3)=5\cdot4\cdot3=60.
	\]
	\end{solution}
\end{question}

\subsection{二项式定理与组合恒等式}

\subsubsection{二项式定理}

\[
(x+y)^n=\sum_{k=0}^{n}\binom nk x^{n-k}y^k.
\]

常用特例：
\[
\sum_{k=0}^{n}\binom nk=2^n,
\]
\[
\sum_{k=0}^{n}(-1)^k\binom nk=0\quad(n\ge1).
\]

\subsubsection{常用组合恒等式}

\[
\binom nr=\binom n{n-r},
\]
\[
\binom nr=\binom{n-1}r+\binom{n-1}{r-1},
\]
\[
\sum_{k=0}^{r}\binom mk\binom n{r-k}=\binom{m+n}r.
\]

\subsection{鸽巢原理}

\subsubsection{基本形式}

若把 \(n+1\) 个对象放入 \(n\) 个盒子，则至少有一个盒子中含有不少于 \(2\) 个对象。

\subsubsection{推广形式}

若把 \(N\) 个对象放入 \(k\) 个盒子，则至少有一个盒子中含有不少于
\[
\left\lceil\frac Nk\right\rceil
\]
个对象。

\begin{question}{例题2}
	从 \(13\) 个人中证明必有两个人出生月份相同。
	
	\begin{solution}
	一年有 \(12\) 个月，把 \(13\) 个人按出生月份放入 \(12\) 个盒子。由鸽巢原理，至少有一个月份中有不少于
	\[
	\left\lceil\frac{13}{12}\right\rceil=2
	\]
	个人。
	\end{solution}
\end{question}

\subsection{容斥原理的计数应用}

设有限全集为 \(E\)，性质集合为 \(A_1,A_2,\dots,A_n\)。满足至少一个性质的对象数为
\[
\left|\bigcup_{i=1}^{n}A_i\right|.
\]
没有任何性质的对象数为
\[
|E|-\left|\bigcup_{i=1}^{n}A_i\right|.
\]

\begin{question}{例题3}
	求 \(1\) 到 \(100\) 中能被 \(2\) 或 \(3\) 整除的整数个数。
	
	\begin{solution}
	设 \(A\) 为能被 \(2\) 整除的数，\(B\) 为能被 \(3\) 整除的数。
	\[
	|A|=\left\lfloor\frac{100}{2}\right\rfloor=50,\quad
	|B|=\left\lfloor\frac{100}{3}\right\rfloor=33,
	\]
	\[
	|A\cap B|=\left\lfloor\frac{100}{6}\right\rfloor=16.
	\]
	所以
	\[
	|A\cup B|=50+33-16=67.
	\]
	\end{solution}
\end{question}

\subsection{递推关系}

\subsubsection{递推关系与初始条件}

用前若干项表示后一项的关系称为递推关系。递推关系加上足够的初始条件才能确定唯一数列。

例如斐波那契数列：
\[
F_0=0,\quad F_1=1,\quad F_n=F_{n-1}+F_{n-2}\quad(n\ge2).
\]

\subsubsection{一阶线性递推}

若
\[
a_n=ra_{n-1}+c,
\]
其中 \(r,c\) 为常数，则当 \(r\ne1\) 时
\[
a_n=r^na_0+c\frac{r^n-1}{r-1}.
\]
当 \(r=1\) 时
\[
a_n=a_0+nc.
\]

\begin{question}{例题4}
	求递推
	\[
	a_n=2a_{n-1}+1,\quad a_0=0
	\]
	的通项。
	
	\begin{solution}
	套用一阶线性递推公式，\(r=2,c=1\)，得
	\[
	a_n=2^n\cdot0+\frac{2^n-1}{2-1}=2^n-1.
	\]
	\end{solution}
\end{question}

\subsubsection{二阶常系数齐次线性递推}

设
\[
a_n=c_1a_{n-1}+c_2a_{n-2}.
\]
其特征方程为
\[
\lambda^2-c_1\lambda-c_2=0.
\]

\begin{enumerate}
	\item 若有两个不同实根 \(\lambda_1,\lambda_2\)，则
	\[
	a_n=A\lambda_1^n+B\lambda_2^n.
	\]
	\item 若有重根 \(\lambda\)，则
	\[
	a_n=(A+Bn)\lambda^n.
	\]
\end{enumerate}

常数 \(A,B\) 由初始条件确定。

\begin{question}{例题5}
	求
	\[
	a_n=3a_{n-1}-2a_{n-2},\quad a_0=1,\ a_1=2
	\]
	的通项。
	
	\begin{solution}
	特征方程为
	\[
	\lambda^2-3\lambda+2=0,
	\]
	根为 \(\lambda_1=1,\lambda_2=2\)。所以
	\[
	a_n=A+B2^n.
	\]
	由 \(a_0=1\) 得 \(A+B=1\)，由 \(a_1=2\) 得 \(A+2B=2\)，解得 \(B=1,A=0\)。因此
	\[
	a_n=2^n.
	\]
	\end{solution}
\end{question}

\subsection{生成函数简介}

数列 \(\{a_n\}\) 的普通生成函数定义为
\[
G(x)=\sum_{n=0}^{\infty}a_nx^n.
\]
常用生成函数：
\[
\frac{1}{1-x}=1+x+x^2+\cdots,
\]
\[
\frac{1}{(1-x)^2}=\sum_{n=0}^{\infty}(n+1)x^n,
\]
\[
(1+x)^n=\sum_{k=0}^{n}\binom nkx^k.
\]

生成函数常用于求递推通项、组合数列和选取问题的计数。

\subsection{本章重点定义与定理速查}

\subsubsection{计数基本模型}

\begin{enumerate}
	\item \textbf{加法原理}：分类讨论且各类互不重叠时，总数为各类数量之和。
	\item \textbf{乘法原理}：分步完成且每一步选择数确定时，总数为各步选择数之积。
	\item \textbf{排列}：从 \(n\) 个不同元素中取 \(r\) 个排成一列：
	\[
	P(n,r)=\frac{n!}{(n-r)!}.
	\]
	\item \textbf{组合}：从 \(n\) 个不同元素中取 \(r\) 个不考虑顺序：
	\[
	\binom nr=\frac{n!}{r!(n-r)!}.
	\]
	\item \textbf{可重复排列}：从 \(n\) 个元素中每次可重复地取 \(r\) 次并考虑顺序，共 \(n^r\) 种。
	\item \textbf{可重复组合}：从 \(n\) 类元素中取 \(r\) 个，允许重复且不考虑顺序，共
	\[
	\binom{n+r-1}{r}
	\]
	种。这等价于非负整数方程
	\[
	x_1+x_2+\cdots+x_n=r
	\]
	的解数。
	\item \textbf{正整数解计数}：方程
	\[
	x_1+x_2+\cdots+x_n=r,\qquad x_i\ge1
	\]
	的解数为
	\[
	\binom{r-1}{n-1}.
	\]
\end{enumerate}

\subsubsection{二项式与组合恒等式}

\begin{enumerate}
	\item \textbf{二项式定理}：
	\[
	(x+y)^n=\sum_{k=0}^n\binom nkx^{n-k}y^k.
	\]
	\item \textbf{对称公式}：
	\[
	\binom nr=\binom n{n-r}.
	\]
	\item \textbf{Pascal 公式}：
	\[
	\binom nr=\binom{n-1}{r}+\binom{n-1}{r-1}.
	\]
	\item \textbf{Vandermonde 恒等式}：
	\[
	\sum_{k=0}^r\binom mk\binom n{r-k}=\binom{m+n}{r}.
	\]
	\item \textbf{二项式系数和}：
	\[
	\sum_{k=0}^n\binom nk=2^n,\qquad
	\sum_{k=0}^n(-1)^k\binom nk=0.
	\]
	\item \textbf{奇偶项和}：
	\[
	\sum_{\substack{0\le k\le n\\k\text{ 为偶数}}}\binom nk
	=
	\sum_{\substack{0\le k\le n\\k\text{ 为奇数}}}\binom nk
	=2^{n-1}\quad(n\ge1).
	\]
\end{enumerate}

\subsubsection{容斥、鸽巢与错排}

\begin{enumerate}
	\item \textbf{广义鸽巢原理}：把 \(N\) 个对象放入 \(k\) 个盒子，必有一个盒子至少含
	\[
	\left\lceil\frac Nk\right\rceil
	\]
	个对象。
	\item \textbf{抽屉下界模型}：若要保证某盒至少有 \(r\) 个对象，最少需要
	\[
	k(r-1)+1
	\]
	个对象。
	\item \textbf{容斥原理}：
	\[
	\left|\bigcup_{i=1}^nA_i\right|
	=
	\sum_i|A_i|-\sum_{i<j}|A_i\cap A_j|+\sum_{i<j<k}|A_i\cap A_j\cap A_k|-\cdots+(-1)^{n+1}|A_1\cap\cdots\cap A_n|.
	\]
	\item \textbf{错排数}：\(n\) 个元素的排列中没有任何元素在原位置的排列数为
	\[
	D_n=n!\sum_{k=0}^n\frac{(-1)^k}{k!}.
	\]
	递推形式为
	\[
	D_n=(n-1)(D_{n-1}+D_{n-2}),\qquad D_1=0,\ D_2=1.
	\]
\end{enumerate}

\subsubsection{递推与生成函数}

\begin{enumerate}
	\item \textbf{递推关系}：用前面若干项定义后一项的关系。递推关系必须配合足够初始条件才能确定唯一数列。
	\item \textbf{一阶线性递推}：
	\[
	a_n=ra_{n-1}+c.
	\]
	当 \(r\ne1\) 时，
	\[
	a_n=r^na_0+c\frac{r^n-1}{r-1};
	\]
	当 \(r=1\) 时，
	\[
	a_n=a_0+nc.
	\]
	\item \textbf{二阶常系数齐次递推}：
	\[
	a_n=c_1a_{n-1}+c_2a_{n-2}
	\]
	的特征方程是
	\[
	\lambda^2-c_1\lambda-c_2=0.
	\]
	不同根给出 \(A\lambda_1^n+B\lambda_2^n\)，重根给出 \((A+Bn)\lambda^n\)。
	\item \textbf{普通生成函数}：
	\[
	G(x)=\sum_{n=0}^{\infty}a_nx^n.
	\]
	它把数列问题转化为形式幂级数的代数运算。
	\item \textbf{常用生成函数}：
	\[
	\frac1{1-x}=\sum_{n=0}^{\infty}x^n,\qquad
	\frac1{(1-x)^2}=\sum_{n=0}^{\infty}(n+1)x^n.
	\]
\end{enumerate}
